
\medskip
\phantomsection\label{prop:sga3-VIB-3.7}%
\noindent\textbf{Proposition 3.7.---}%
\emph{Let \(G\) be an \(S\)-group locally of finite presentation.
Then:}
\begin{enumerate}
\item[\textup{\SGAInlineRefTarget{prop:sga3-VIB-3.7:item:i}(i)}] \emph{\(\underline{G^0}\) is ind-constructible
  in \(G\).}
\item[\textup{\SGAInlineRefTarget{prop:sga3-VIB-3.7:item:ii}(ii)}] \emph{If, moreover, \(G\) is quasi-separated
  over \(S\), and \(S\) is quasi-compact and quasi-separated, then
  \(\underline{G^0}\) is constructible.}
\item[\textup{\SGAInlineRefTarget{prop:sga3-VIB-3.7:item:iii}(iii)}] \emph{Consequently, if \(G\) is
  quasi-separated over \(S\), then \(\underline{G^0}\) is locally
  constructible.}
\end{enumerate}

We first prove the first assertion. Since
\[
  \pi:G\longrightarrow S
\]
is locally of finite presentation, given \(s\in S\), there are an
open subset \(U\) of \(G\) containing \(\epsilon(s)\) and an open
subset \(V\) of \(S\) containing \(s\) such that
\[
  \pi(U)\subset V
\]
and the induced morphism
\[
  \pi':U\longrightarrow V
\]
is of finite presentation. Then
\[
  T=\epsilon^{-1}(U)
\]
is an open subset of \(S\). Set
\[
  W=(\pi')^{-1}(T),\qquad
  \pi''=\pi'|_W.
\]
Thus \(\pi'':W\to T\) is of finite presentation and admits as a
section the morphism
\[
  \epsilon'':T\longrightarrow W
\]
induced by \(\epsilon\).

For every \(t\in T\), since \(\underline{G_t^0}\) is irreducible
(VI\(_A\), \hyperref[prop:sga3-VIA-2.4]{2.4}), \(W\cap\underline{G_t^0}\) is dense in
\(\underline{G_t^0}\), hence irreducible and therefore connected.
It is thus the connected
component of \((\pi'')^{-1}(t)\) containing \(\epsilon''(t)\).
It follows from EGA~IV\(_3\), 9.7.12, that the union \(W^0\) of the
\(W\cap\underline{G_t^0}\), for \(t\in T\), is locally
constructible in \(W\).
On the other hand, VI\(_A\), \SGARefLink{sga3:vi:proposition:0-5}{0.5}, gives
\[
  \underline{G^0}\cap\pi^{-1}(T)=W^0\cdot W^0;
\]
that is, \(\underline{G^0}\cap\pi^{-1}(T)\) is the image of
\[
  W^0\times_T W^0
\]
under the morphism
\[
  \mu'':W\times_T W\longrightarrow\pi^{-1}(T)
\]
induced by \(\mu\).%
\footnote{\label{note:sga3-VIB-25}Editor's note: we have expanded the
original in what follows.}
Since \(W\times_T W\) (respectively, \(\pi^{-1}(T)\)) is of finite
presentation (respectively, locally of finite presentation) over
\(T\), the morphism \(\mu''\) is locally of finite presentation and
quasi-separated, by EGA~IV\(_1\), 1.4.3 and 1.2.2. If, moreover,
\(\pi^{-1}(T)\) is quasi-separated over \(T\), then \(\mu''\) is
quasi-compact (loc.\ cit., 1.2.4), hence of finite presentation.
Since \(W^0\times_T W^0\) is locally constructible in \(W\times_T W\)
(because \(W^0\) is locally constructible in \(W\)), Chevalley's
constructibility theorem (loc.\ cit., 1.8.4 and 1.9.5(viii)) shows
that \(\underline{G^0}\cap\pi^{-1}(T)\) is ind-constructible in
\(\pi^{-1}(T)\), and is locally constructible in \(\pi^{-1}(T)\)
if \(G\), and hence \(\pi^{-1}(T)\), is quasi-separated over \(T\).
This proves assertions \textup{\SGARefLink{prop:sga3-VIB-3.7:item:i}{(i)}} and \textup{\SGARefLink{prop:sga3-VIB-3.7:item:iii}{(iii)}}.

Now suppose that \(G\) is quasi-separated over \(S\), and that \(S\)
is quasi-compact and quasi-separated. By \hyperref[prop:sga3-VIB-3.5]{3.5}, there is a
quasi-compact open subset \(U\) of \(G\) containing
\(\underline{G^0}\). Since
\(G\) is quasi-separated over \(S\), \(G\) is quasi-separated;
hence the open subset \(U\) is retrocompact (EGA~IV\(_1\), 1.2.7),
and it suffices to show that \(\underline{G^0}\) is constructible
in \(U\)
(EGA~0\(_{\mathrm{III}}\), 9.1.8).

Moreover, \(U\) is quasi-compact, hence quasi-compact over \(S\)
(EGA~IV\(_1\), 1.2.4), and is quasi-separated over \(S\).
Therefore the restriction of \(\pi\) to \(U\) is of finite
presentation. By EGA~IV\(_3\), 9.7.12, \(\underline{G^0}\) is locally
constructible in \(U\), hence constructible in \(U\), since \(U\)
is quasi-compact and quasi-separated (EGA~IV\(_1\), 1.8.1).
This proves \textup{(ii)}, and it follows that, for every
quasi-compact and quasi-separated open subset \(T\) of \(S\)
(for example, every affine open subset of \(S\)),
\[
  \underline{G^0}\cap\pi^{-1}(T)
\]
is constructible.

\medskip
\phantomsection\label{cor:sga3-VIB-3.8}%
\noindent\textbf{Corollary 3.8.---}
\emph{Let \(S_0\) be a quasi-compact and quasi-separated scheme,
\(I\) an increasing filtered preordered set,
\((\mathcal A_i)_{i\in I}\) a direct system of commutative
quasi-coherent \(\mathcal O_{S_0}\)-algebras,}
\[
  \mathcal A=\varinjlim_i\mathcal A_i,\qquad
  S_i=\operatorname{Spec}\mathcal A_i\quad(i\in I),\qquad
  S=\operatorname{Spec}\mathcal A
\]
\emph{\textup{(cf.\ EGA~II, 1.3.1)}. Let \(G\) be an
\(S_0\)-group scheme locally of finite presentation. Then the
canonical map}
\[
  \varinjlim_i G^0(S_i)\longrightarrow G^0(S)
\]
\emph{is bijective.}

Since \(G\) is locally of finite presentation over \(S_0\),%
\begingroup
\renewcommand{\thefootnote}{\fnsymbol{footnote}}%
\footnote[1]{\label{note:sga3-VIB-3.8-base}Translator's note: the
French source prints ``over \(S\)'' here. The scheme \(G\) was
defined over \(S_0\), while \(S\) and the \(S_i\) are the test
schemes of the direct system; the required base is therefore
\(S_0\).}%
\endgroup{}
the canonical map
\[
  \varinjlim_i G(S_i)\longrightarrow G(S)
\]
is bijective, by EGA~IV\(_2\), 8.14.2(c). It follows immediately
that the canonical map
\[
  \varinjlim_i G^0(S_i)\longrightarrow G^0(S)
\]
is injective. We prove that it is surjective. Let
\[
  g\in G^0(S)\subset G(S).
\]
There is an \(i\in I\) such that \(g\) factors, by means of an
element \(g_i\in G(S_i)\), through \(S\to S_i\). By hypothesis,
\[
  g^{-1}(\underline{G^0})=S.
\]
But \(\underline{G^0}\) is ind-constructible in \(G\), by \hyperref[prop:sga3-VIB-3.7]{3.7}, so
\[
  g_i^{-1}(\underline{G^0})
\]
is ind-constructible in \(S_i\). It follows from EGA~IV\(_2\),
8.3.4, that there is an index \(j\geq i\) such that
\[
  g_j^{-1}(\underline{G^0})=S_j,
\]
where \(g_j\) is the map induced by \(g_i\) under the base change
\(S_j\to S_i\). Thus \(g_j\in G^0(S_j)\), and hence \(g\) comes
from an element of
\[
  \varinjlim_i G^0(S_i).
\]
